\subsection{Pipeline 3: Moondream Prototyping to Local CNN-Based Hybrid}

At the start of this track, we first experimented with vision-language models (mainly Moondream and LLaVA) to speed up labeling. In practice, API-key based inference had provider limits, quota/rate constraints, and unstable throughput. For a dataset size of 10,000 images, relying on external API calls was difficult to scale reliably.

Therefore, we moved to a local reproducible pipeline. The finalized method in this section is a \textbf{CNN-based hybrid transfer/fine-tuning workflow} (CNN backbone with practical preprocessing and shape-prior adaptation), rather than a purely standalone end-to-end CNN-only setup.

\subsubsection{Pipeline 3 Stages (Final Local Workflow)}

\begin{enumerate}
    \item \textbf{Initial staged training of the local CNN-based hybrid model on \texttt{AHDBase\_TrainingSet}.}
    The script \path{train_digits_staged.py} produced \path{runs/digits_staged}. The summary file reports best validation accuracy of \textbf{0.98125}.

    \item \textbf{Transfer evaluation on \texttt{PracticalGroundTruth500}.}
    The script \path{predict_digits_cnn.py} was used to compare transfer checkpoints; \path{stage_05.pt} transferred better than \path{best_model.pt} (43.00\% vs 41.20\%).

    \item \textbf{Support/query prototype adaptation on practical subset.}
    Using 3 support images per class (seed 7), autocontrast, and shape prior, the measured query accuracy reached \textbf{413/470 = 87.87\%}.

    \item \textbf{Fine-tuning on \texttt{PracticalGroundTruth500}.}
    Starting from \path{runs/digits_staged/stage_05.pt}, fine-tuning produced \path{runs/practical500_finetune/best_model.pt} with:
    \begin{itemize}
        \item best validation accuracy = \textbf{0.94}
        \item held-out test accuracy = \textbf{0.84}
    \end{itemize}

    \item \textbf{Full prediction on \texttt{Indian\_Digits\_Train}.}
    The final full run used \path{runs/practical500_finetune/best_model.pt} and produced predictions for all \textbf{10,000} images with no empty \texttt{predicted\_label} values.

    \item \textbf{Submission ordering correction.}
    Final predictions were sorted by numeric image ID (1,2,3,...) to avoid lexicographic-order mismatch during oracle-style evaluation.
\end{enumerate}

\subsubsection{Key Commands Used}

\noindent\textbf{Fine-tuning stage command.}
\begin{lstlisting}[language=bash]
python .\train_digits_staged.py \
    --data-dir .\PracticalGroundTruth500 \
    --resume-from .\runs\digits_staged\stage_05.pt \
    --stage-sizes 5,10,20,30,40 \
    --epochs-per-stage 20 \
    --patience 5 \
    --lr 1e-4 \
    --val-ratio 0.10 \
    --test-ratio 0.10 \
    --output-dir .\runs\practical500_finetune
\end{lstlisting}

\noindent\textbf{Final full prediction command.}
\begin{lstlisting}[language=bash]
python .\predict_digits_cnn.py \
    --checkpoint .\runs\practical500_finetune\best_model.pt \
    --input-dir .\Indian_Digits_Train \
    --support-dir .\PracticalGroundTruth500 \
    --autocontrast \
    --use-shape-prior \
    --shape-prior-weight 0.3 \
    --output-csv .\runs\indian_digits_cnn_finetuned_predictions.csv
\end{lstlisting}

\subsubsection{MATLAB Online Oracle Scoring (\texttt{check\_accuracy.p})}

For final local-model validation, we used MATLAB Online and the provided P-coded oracle file \path{check_accuracy.p}. The scoring input was a CSV prediction file containing labels for all 10,000 images.

\paragraph{MATLAB checker used for \texttt{predictions\_final.csv}.}
\begin{lstlisting}[language=Matlab]
% reduced MATLAB Online checker for predictions_final.csv
pred_file = 'predictions_final.csv';
T = readtable(pred_file);
fprintf('Scoring %s\n', pred_file);

assert(ismember('predicted_label', T.Properties.VariableNames), ...
    'predicted_label column not found');

preds = str2double(string(T.predicted_label));  % handles numeric/text columns
assert(~any(isnan(preds)), ...
    'Found %d missing predicted_label entries.', sum(isnan(preds)));
assert(numel(preds) == 10000, ...
    'predicted vector must have length 10000');

% Align to oracle order (1..10000) before scoring.
n = numel(preds);
ord = (1:n)';

if ismember('image_id', T.Properties.VariableNames)
    ids = str2double(string(T.image_id));
    assert(~any(isnan(ids)), ...
        'Found %d invalid image_id entries.', sum(isnan(ids)));
    [ids_sorted, ord] = sort(ids(:));
    assert(isequal(ids_sorted, (1:n)'), ...
        'image_id must be exactly 1..%d with no duplicates/missing rows.', n);
    preds = preds(ord);
    fprintf('Aligned predictions using image_id.\n');
elseif ismember('path', T.Properties.VariableNames)
    p = string(T.path);
    ids = nan(n, 1);
    for i = 1:n
        tok = regexp(p(i), '(\\|/)(\d+)\.[^\\/]+$', 'tokens', 'once');
        if ~isempty(tok)
            ids(i) = str2double(tok{2});
        end
    end
    assert(~any(isnan(ids)), ...
        ['Could not parse numeric image id from %d path entries. ' ...
         'Expected .../<number>.ext'], sum(isnan(ids)));
    [ids_sorted, ord] = sort(ids(:));
    assert(isequal(ids_sorted, (1:n)'), ...
        'Parsed path IDs must be exactly 1..%d with no duplicates/missing rows.', n);
    preds = preds(ord);
    fprintf('Aligned predictions using numeric ID parsed from path.\n');
else
    fprintf(['No image_id/path column found; using row order as-is. ' ...
             'Ensure rows already match oracle order.\n']);
end

preds_int = int32(preds(:));
try
    [acc, n_correct, n_total] = check_accuracy(preds_int);
    fprintf('accuracy=%.6f\ncorrect=%d/%d\n', acc, n_correct, n_total);
catch
    acc = check_accuracy(preds_int);
    fprintf('accuracy=%g\n', acc);
end
\end{lstlisting}

\paragraph{Output snippet (MATLAB Online).}
\begin{lstlisting}[language=Matlab]
>> check
Scoring predictions_final.csv
Aligned predictions using numeric ID parsed from path.

===== Classification Results =====
  Correct : 9339 / 10000
  Accuracy: 93.39%
==================================
\end{lstlisting}

\paragraph{Consistency note across Part 3 pipelines.}
The 93.39\% value above is the oracle score for the specific file \texttt{predictions\_final.csv} in this CNN run. It is reported independently from Pipeline 1 (99.09\%) and Pipeline 2 (99.07\%) results, since each pipeline uses different training/adaptation procedures and run configurations.

\subsubsection{Final Artifacts and Results}

\begin{table}[H]
\centering
\caption{Pipeline 3 final artifacts and measured outcomes}
\label{tab:part3_pipeline3_summary}
\begin{tabular}{>{\raggedright\arraybackslash}p{0.46\textwidth}>{\raggedright\arraybackslash}p{0.46\textwidth}}
\toprule
\textbf{Item} & \textbf{Value} \\
\midrule
Final checkpoint & \path{runs/practical500_finetune/best_model.pt} \\
Full prediction file & \path{runs/indian_digits_cnn_finetuned_predictions.csv} \\
Numeric-sorted prediction file & \path{runs/indian_digits_cnn_finetuned_predictions_numeric.csv} \\
Recommended submission file & \path{runs/indian_digits_cnn_finetuned_submission.csv} \\
Labels-only alternative & \path{runs/indian_digits_cnn_finetuned_labels_only.txt} \\
Practical subset match check (diagnostic) & 470/500 = 94.0\% \\
Oracle score for shown \texttt{predictions\_final.csv} & 9339/10000 = 93.39\% \\
Predictions generated in final run & 10,000 / 10,000 \\
\bottomrule
\end{tabular}
\end{table}

\paragraph{Manual effort note.}
For the 10,000-image final inference stage, no additional per-image manual labeling was needed; the pipeline relied on the already available labeled practical subset (\path{PracticalGroundTruth500}) plus local model inference.

\subsubsection{Additional Pipeline 3 Baseline: Qwen2-VL (Practical 500 Set)}

After finalizing the local CNN-based hybrid track, we also evaluated a separate Qwen2-VL baseline as an additional reference experiment. This Qwen run focused on the practical set workflow (500 labeled samples for evaluation) and was used to measure how far prompt/inference-only tuning can go without task-specific model fine-tuning.

\paragraph{Qwen baseline setup.}
\begin{itemize}
    \item Task: classify handwritten Arabic-Indic digits (0--9).
    \item Evaluation set: \path{Qwen/PracticalGroundTruth500} (500 labeled samples, class-balanced with 50 samples per class).
    \item Main workflow notebooks: \path{Qwen/scripts/create_labels_and_predection.ipynb}, \path{Qwen/scripts/resize_images.ipynb}, and \path{Qwen/scripts/compute_accuracy.ipynb}.
    \item Key hardening steps: robust output parsing, stable ground-truth merge, preprocessing with aspect-ratio-preserving resize/padding, and GPU-enabled execution path.
\end{itemize}

\paragraph{Qwen results summary.}
\begin{table}[H]
\centering
\caption{Pipeline 3 Qwen2-VL practical-set results}
\label{tab:part3_qwen_summary}
\begin{tabular}{>{\raggedright\arraybackslash}p{0.48\textwidth}>{\raggedright\arraybackslash}p{0.44\textwidth}}
\toprule
\textbf{Item} & \textbf{Value} \\
\midrule
Final practical accuracy & 112/500 = \textbf{22.40\%} \\
Invalid predictions in final run & \textbf{0} \\
Early failed state & 0/500 = 0.0\% (dominated by invalid parsing) \\
Intermediate stable checkpoints & 73/500 = 14.6\%, then 111/500 = 22.2\% \\
Prediction distribution note (final) & Strong class bias remained (e.g., class 2 over-predicted) \\
\bottomrule
\end{tabular}
\end{table}

\paragraph{Assignment reporting note.}
This Qwen2-VL result is reported as an additional practical-set baseline inside Pipeline 3. Since this run was evaluated on the practical 500-sample set (not full oracle scoring on 10,000 predictions), it is used for methodology comparison and debugging evidence, while the official Pipeline 3 full-dataset oracle figure remains the local-model result (93.39\%).

\subsection{Part 3 Cross-Pipeline Comparison}

Table~\ref{tab:part3_cross_pipeline} consolidates the finalized outcomes, including supplementary add-on baselines, so all reported Part 3 results can be compared directly.

\begin{table}[H]
\centering
\caption{Part 3 cross-pipeline comparison}
\label{tab:part3_cross_pipeline}
\begin{tabular}{>{\raggedright\arraybackslash}p{0.26\textwidth}>{\raggedright\arraybackslash}p{0.14\textwidth}>{\raggedright\arraybackslash}p{0.10\textwidth}>{\raggedright\arraybackslash}p{0.20\textwidth}>{\raggedright\arraybackslash}p{0.24\textwidth}}
\toprule
\textbf{Pipeline} & \textbf{Oracle Accuracy} & \textbf{Iterations} & \textbf{Manual Effort} & \textbf{Observation} \\
\midrule
Pipeline 1 (K-means bootstrapping + HITL SVM) & 99.09\% & 3 & 424 handled images, 1240.0 s & Highest oracle accuracy with the lowest estimated manual time among the two HITL-SVM pipelines. \\ \\
Pipeline 2 (active learning + pseudo-labeling SVM) & 99.07\% & 8 & 460 labels, 4600 s & Reached the 99\% target with stable gains, but required more manual actions and time than Pipeline 1. \\ \\
Pipeline 3 (Moondream prototyping -> local CNN-based hybrid) & 93.39\% & N/A & No additional per-image labels during final inference; reused \path{PracticalGroundTruth500} & Started from VLM prototyping, then shifted to a local CNN-based hybrid baseline; lower oracle accuracy under this run setup. \\ \\
Pipeline 3 add-on (Qwen2-VL prompt/inference baseline) & N/A (practical: 22.40\%) & N/A & No new labels; evaluated using existing 500 practical labels & Stable end-to-end after debugging (invalid=0), but strong class-bias limited final practical accuracy. \\ \\
\bottomrule
\end{tabular}
\end{table}

\paragraph{Comparative interpretation.}
\begin{itemize}
    \item Pipeline 1 and Pipeline 2 are nearly tied in oracle accuracy (difference = 0.02 percentage points).
    \item Pipeline 1 achieved this with less estimated manual effort than Pipeline 2 in the reported runs.
    \item Pipeline 3's local CNN-based hybrid run (93.39\% oracle) substantially outperformed the Qwen practical-set baseline (22.40\% practical).
    \item The Qwen baseline demonstrates a working VLM path, but practical-only scoring and class-bias behavior make it a supplementary reference, not the primary Pipeline 3 result.
\end{itemize}